% Sets and expectation
\newcommand{\R}{\mathbb{R}}
\newcommand{\N}{\mathbb{N}}
\newcommand{\Exp}{\mathbb{E}} % use \Exp for expectation to avoid clashing with edge set E

% Calligraphic sets
\newcommand{\cS}{\mathcal{S}}
\newcommand{\cT}{\mathcal{T}}     % set of spanning trees
\newcommand{\cF}{\mathcal{F}}
\newcommand{\Scenarios}{\mathcal{U}} % uncertainty set

% Operators and helpers
\newcommand{\OPT}{\mathrm{OPT}}
\newcommand{\ALG}{\mathrm{ALG}}
\newcommand{\opt}{\OPT} % backwards-compatible alias

\DeclareMathOperator{\cost}{cost}
\DeclareMathOperator{\regret}{Regret}

\newcommand{\Regret}[2]{\regret(#1,#2)} % usage: \Regret{T}{c} gives Regret(T,c)

\newcommand{\norm}[1]{\left\lVert #1 \right\rVert}
\newcommand{\abs}[1]{\left\lvert #1 \right\rvert}
\newcommand{\st}{\text{s.t.}}

\newcommand{\MST}{\mathrm{MST}}
\newcommand{\MSTcost}[1]{\mathrm{MST}(#1)}
\newcommand{\cs}[1]{c^{(#1)}}
\newcommand{\wc}[1]{\mathrm{wc}(#1)} % worst-case cost: usage \wc{T} gives wc(T) = max_{c in U} c(T)
\newcommand{\wcr}[1]{\mathrm{wcr}(#1)} % maximum regret: usage \wcr{T} gives wcr(T) = max_{c in U} Regret(T,c). NB \wr is LaTeX's wreath product and cannot be used.

% Graph notation
\newcommand{\cut}[1]{\delta(#1)} % cut notation, usage: \cut{X}

% Complexity helpers
\DeclareMathOperator{\poly}{poly} % polynomial, usage: O(\poly(n))